\section{Reproducibility}
\label{sec:reproducibility}

The code, the committed forecasts and the scripts are in the repository named in the header block, under \texttt{manuscripts/adaptive-conformal-recalibration-of-var}. The results were produced with release v1.00.000 of CAOS\_Contraste (tag v1.00.000). The engine is \texttt{riskvalidation} 1.1.0, released on PyPI, with the extras for market risk, machine learning and stress tests, and the interpreter is Python 3.13 with \texttt{numpy} 2.4, \texttt{scipy} 1.18, \texttt{arch} 8.0 and \texttt{torch} 2.14 as pinned by the pipeline. The committed documents of the market case under \texttt{data/derived/C11} hold the forecasts of the ladder and the wrapped runs of the three models of the hypothesis. The raw series (the French library file, the Treasury par curves and the H.10 release) are read from the data root named by the environment variable \texttt{CONTRASTE\_DATA}; the French file is distributed for derived use only. % voice-ok: citation of the product release the results were produced with

Every number of the text is generated. The script \texttt{make\_numbers.py} writes \texttt{tables/numbers.tex}, which defines a macro for each quoted figure; \texttt{make\_tables.py} and \texttt{make\_figures.py} write the tables and the vector figures from the JSON outputs of the experiments, which are committed under \texttt{experiments/}. The experiments are run from the repository root:
\begin{description}
\item[\texttt{e2\_wrap\_ladder.py}] the wrapper on ten ladder models in six windows, with the reproduction of the committed rows (\reproEqual{} of \reproChecked{} rows equal to the committed documents, and the unwrapped runs of all ten models equal to the ladder's forecasts day by day); \texttt{python manuscripts/adaptive-conformal-recalibration-of-var/experiments/e2\_wrap\_ladder.py --workers 28 --cache \$CAOS\_TEMP\_ROOT/contraste-adaptive-conformal/cache}, with \texttt{CONTRASTE\_DATA} set;
\item[\texttt{e3\_ablation.py}] the ablation of \cref{tab:ablation}; \texttt{python manuscripts/adaptive-conformal-recalibration-of-var/experiments/e3\_ablation.py --workers 28 --returns-cache \$CAOS\_TEMP\_ROOT/contraste-adaptive-conformal/windows.npz}, with \texttt{CONTRASTE\_DATA} set;
\item[\texttt{e4\_score\_space.py}] the two score-space recalibrations on all eighteen models; \texttt{python manuscripts/adaptive-conformal-recalibration-of-var/experiments/e4\_score\_space.py --returns-cache \$CAOS\_TEMP\_ROOT/contraste-adaptive-conformal/windows.npz}, with \texttt{CONTRASTE\_DATA} set;
\item[\texttt{check\_e4\_baseline.py}] the consistency check of the previous output: the unwrapped baseline of experiment 4 equals that of experiment 2 and the committed documents; \texttt{python manuscripts/adaptive-conformal-recalibration-of-var/experiments/check\_e4\_baseline.py};
\item[\texttt{e5\_generators.py}] the generators of known truth, 100 replications of four scenarios, the seed of replication $k$ of scenario $s$ being \seedBase{} plus $1000s + k$; \texttt{python manuscripts/adaptive-conformal-recalibration-of-var/experiments/e5\_generators.py --workers 28};
\item[\texttt{e6\_theory.py}] the numerical checks of the propositions: the identity on the 18 real level paths, the clustered sequences, the scale sweep with 50 replications of 5,000 days (seed \seedBase{} plus 777000 plus the replication) and the ES ratios; \texttt{python manuscripts/adaptive-conformal-recalibration-of-var/experiments/e6\_theory.py};
\item[\texttt{e7\_dispersion.py}] the dispersion of the hit count and the size of the tests on wrapped runs (\cref{sec:results-dispersion}); \texttt{python manuscripts/adaptive-conformal-recalibration-of-var/experiments/e7\_dispersion.py}; it reads \texttt{e2\_wrap\_ladder.json} and \texttt{e5\_generators.json}, and its own draws use \texttt{default\_rng} with the seed list \seedBase{}, 7 or 8, and the indices of the row.
\end{description}
The scripts \texttt{make\_numbers.py}, \texttt{make\_tables.py} and \texttt{make\_figures.py} are run from the repository root without arguments after the experiments.
The bootstrap uses circular blocks of 10 days and 2,000 draws with the seed \seedBase{} plus a fixed integer per window, model and level. The cost of the experiments was measured on a workstation with 32 logical cores: the wrapped runs of ten models in six windows (the heavy part is the daily refit of the GARCH, GJR, EGARCH and extreme-value models) took 168 seconds of wall time with 28 worker processes from an empty run cache. The ablation took 439 seconds, the score-space recalibrations 130 seconds and the generators 176 seconds in a first run, and 90, 16 and 31 seconds in a repeat that reused the cached parsed returns and fits, the times each script records in its output; the numerical checks of the theory took 613 seconds, measured outside the script. A repeat of experiments 2, 3, 5 and 6 reproduces the stored outputs. The score-space recalibrations take 32 seconds, and \texttt{check\_e4\_baseline.py} verifies that each comparator run scores the pair of its own VaR and ES; the dispersion experiment takes under two minutes. Each script finishes in under 30 minutes. The step of the wrapper, 0.005, and the band, $[0.001, 0.25]$, are the defaults of the engine and appear nowhere else as parameters.

The numerical checks of the theory are part of these scripts. The identity of \cref{prop:identity} holds on the 18 real level paths with a largest residual of $\thAResidual$ and the bound holds on \thAHold{} of \thA{}; the clustered sequences of \cref{tab:cluster} are tested with the engine's own tests; the scale sweep of \cref{tab:scale-sweep} is a simulation with the stated seeds; and the ES ratios of \cref{prop:es} are computed in closed form from the Student-$t$ law. A reader who runs the figure and table scripts from the committed JSON obtains the same files without the raw series; a reader who reruns the experiments needs the series and obtains the same numbers to the last digit of the stored JSON, because no step draws a random number except through the stated seeds.
